\MFQTransition{持仓与收益计算与应用分析答案}

\begin{enumerate}[leftmargin=*]
  \item purge 删除训练标签终点跨入验证区的样本；embargo 在验证区后留出空档，阻断持仓尚未结束或相邻标签带来的依赖。画出训练、验证、测试区间并标注被删行，才能证明两条信息路径都被截断。
  \item 目标仓位是意图，实际仓位由上一期仓位加实际成交增量得到：$p_t=p_{t-1}+\mathrm{fill}_t$。部分成交、停牌、借券失败和延迟都使 $p_t\ne p_t^{target}$；收益必须用 $p_t$ 乘下一期回报，不能用目标直接计算。
  \item 初始仓位 $[1,0,0,0]$，第二期目标反向，订单为 $-2$，成交一半即卖出 1 单位，仓位变为 $[0,0,0,0]$。毛收益 $1\times0.02=0.02$，换手为开仓 1 加平仓 1 共 2，成本 $2\times0.002=0.004$，净收益 $0.016$；每行都保留时间、fill 和现金。
  \item 持有期改为 2 时，新信号不能立即翻仓；要先确定持仓年龄是成交后加一还是决策前加一，再把年龄状态写入递推。逐期表列 target、order、fill、position 和 turnover，比较 $H=1$ 与 $H=2$ 的重叠和成本。
  \item 每笔订单保存可借数量、成交比例、失败原因和时间；失败 fill=0 时旧仓位保持不变，成本也为零但机会成本可单独记录。不能删除失败行，否则剩余样本条件化在可成交事件上。
  \item 全样本标准化使用了决策后的分布信息，应在当时可用的数据上估计变换。未成交订单无法改变实际持仓，收益应按旧持仓加实际成交增量计算。成交比例属于 $[0,1]$；停牌或借券限制还会改变允许的成交量。
  \item 综合项目可以从价差模型的预测出发，逐期计算目标仓位、实际成交、持仓与现金。与简单基线比较时保持相同数据和成本定义，再改变关系稳定性、成交比例或持有期，解释净收益的变化。
\end{enumerate}

\subsection*{基础衔接：独立计算}
净损益为 $1-0.2=0.8$。之后的估计不会改写本期持仓；它只影响下一次调仓及其现金流。
